%%%%%%%%%%  scan idx 618  |  folio 605  %%%%%%%%%%
\[
  R^2f_*(I\otimes f^*M)\simeq \underline{\Hom}_{{O}_S}(Q,M).
\]
\item \uline{Le
conoyau} \(E\) \uline{de} \(f_*({O}_{X'})\to R^1f_*(I)\) \uline{est un} \({O}_S\)\uline{-module localement libre.}
\item \uline{L'homomorphisme} \((f_*{O}_X)_{\red}\longrightarrow (f_*{O}_{X'})_{\red}\) \uline{est surjectif.}
\end{enumerate}
\uline{Alors le morphisme canonique} \(\underline{\Pic}_{X/S}\longrightarrow \underline{\Pic}_{X'/S}\) \uline{est représentable par un morphisme affine de présentation finie~; plus précisément, on a une suite exacte canonique de faisceaux} f.p.p.f.,
\[
  0\longrightarrow \mathbb V(E^\vee)\longrightarrow \underline{\Pic}_{X/S}
  \longrightarrow\underline{\Pic}_{X'/S}\xrightarrow{u}\mathbb V(Q)
\]
\uline{et, si} \(P=\Ker(u)\)\uline{, le monomorphisme} \(P\hookrightarrow\Pic_{X'/S}\) \uline{est représentable par une immersion fermée et} \(\Pic_{X/S}\) \uline{est un torseur sur} \(P\) \uline{sous le groupe} \(V(\check E)\) \uline{(c'est-à-dire, la suite de faisceaux} f.p.p.f.
\[
  0\longrightarrow \mathbb V(\check E)\longrightarrow \Pic_{X/S}\longrightarrow P\longrightarrow0
\]
\uline{est exacte).}

Considérons pour un \(S\)-schéma \(T\) les deux diagrammes suivants à lignes exactes :
\[
\begin{tikzcd}
0 \arrow[r] & I_T \arrow[r] & {O}_{X_T} \arrow[r,dash] & {O}_{X'_T} \arrow[r] & 0\\
0 \arrow[r] & I_T \arrow[r,"\alpha",dash] \arrow[u,equal] & {O}_{X_T}^{*} \arrow[r] \arrow[u,hook] & {O}_{X'_T}^{*} \arrow[r] \arrow[u,hook] & 0
\end{tikzcd}
\qquad (\alpha(i)=1+i)
\]
\[
\begin{tikzcd}[column sep=tiny]
(f_T)_*({O}_{X'_T}) \arrow[r,"\delta_T"] &
R^1(f_T)_*(I_T) \arrow[d,equal]\\
(f_T)_*({O}_{X'_T}^{*}) \arrow[r,"\delta_T^*"] \arrow[u] &
R^1(f_T)_*(I_T) \arrow[r] &
R^1(f_T)_*({O}_{X_T}^{*}) \arrow[r] &
R'(f_T)_*({O}_{X'_T}^{*}) \arrow[r] &
R^2(f_T)_*(I_T) .
\end{tikzcd}
\]

On fera attention que le carré de gauche de ce diagramme \uline{n'est pas nécessairement commutatif}. Cependant on a~:
\[
  \textbf{3.2.}\qquad \im\delta_T^*=\im\delta_T.
\]
